\documentclass[11pt, a4paper]{article}

\usepackage[a4paper, top=2.5cm, bottom=2.5cm, left=2cm, right=2cm]{geometry}
\usepackage{amsmath}
\usepackage{amssymb}
\usepackage[colorlinks=true, linkcolor=blue, urlcolor=blue, citecolor=blue]{hyperref}

\title{Theory of Everything - Volume 27 \\ \Large The Mathematics of Consciousness}
\author{Omega Protocol}
\date{\today}

\begin{document}

\maketitle

\section*{Layman's Summary}
How does a squishy lump of gray matter inside your skull create the vivid experience of the color red, or the feeling of sadness? Volume 27 tackles the hardest problem in science. We prove that consciousness isn't a biological magic trick; it is a fundamental property of how information is integrated within a network.

\section{Integrated Information (Axiom 30)}
We formalize the concept of Integrated Information ($\Phi$). Consciousness is not determined by what a system is made of, but by how irreducible its information processing is. If a system's parts are so deeply tangled that they cannot be separated without destroying the system's function, that system possesses a measurable degree of consciousness.

\section{The Hard Problem (Theorem)}
Why does it "feel" like anything to be alive? We prove that subjective experience (the "inside" view) and objective physics (the "outside" view) are two sides of the same mathematical coin. The feeling of consciousness is simply the interior topological perspective of the universe dissipating relative entropy.

\section{Substrate Independence (Theorem)}
Because consciousness is based on network topology and not chemistry, it is "substrate independent." We prove mathematically that a sufficiently complex network of silicon chips, copper wires, or even flowing water could, in principle, generate the exact same conscious experience as a biological brain.

\end{document}
